\documentclass[11pt]{article}
\usepackage[margin=25mm]{geometry}
\usepackage{fontspec}
\setmainfont{Latin Modern Roman}
\usepackage{amsmath,amssymb,mathtools}
\usepackage{unicode-math}
\setmathfont{Latin Modern Math}
\usepackage{xurl,hyperref}
\hypersetup{colorlinks=true,urlcolor=blue}
\setcounter{secnumdepth}{0}
\setlength{\parindent}{0pt}
\setlength{\parskip}{5pt}
\setlength{\emergencystretch}{3em}
\newcommand{\sref}[2]{\hyperlink{source:#1:#2}{[S#2]}}
\newcommand{\cataloguescope}[1]{\paragraph{Recorded target.}#1}
\begin{document}
\section{Research Problems in Function Theory — Problem 6.59}\label{statement}

\subsection{Definitions and mathematical statement}

\paragraph{Shared notation.}$\mathbb N=\{1,2,\ldots\}$, and $\mathbb Z,\mathbb Q,\mathbb R,\mathbb C$ denote the integers, rationals, reals and complex numbers. Any other conventions are specified locally.


Let $\widehat{\mathbb C}=\mathbb C\cup\{\infty\}$ be the Riemann sphere, and let $D\subset\widehat{\mathbb C}$ be a domain containing $\infty$. Its complement $K$ is compact. A continuous angle assignment to its connected components means a continuous real-valued function $a:K\to\mathbb R$ constant on each component; equivalently, it is continuous on the quotient space of components. The value $a(K_j)$ specifies an unoriented line direction modulo $\pi$. A rectilinear slit domain has complement consisting of pairwise disjoint closed straight segments, allowing degenerate segments consisting of one point. Boundary components correspond under a conformal equivalence.
\paragraph{Statement recorded in the supplied source.}
Does every such $(D,a)$ admit a conformal equivalence $\phi:D\to\widehat{\mathbb C}\setminus\bigcup_j I_j$ for which the boundary component $K_j$ corresponds to a segment $I_j$ parallel to $e^{ia(K_j)}$? The prescribed directions are allowed to have infinitely many distinct values. \sref{statement}{2}
Update 6.59 records FitzGerald and Weening's existence result when the angle assignment has finite range. \sref{statement}{2}
\subsection{Short English statement}
Can an arbitrary continuously prescribed direction for each complementary component be realized by a conformal straight-slit representation?
\subsection{Sources}
\begin{description}
\item[\hypertarget{source:statement:1}{S1}] ULAM.AI and the UnsolvedMath Contributors, \emph{UnsolvedMath: A Curated Collection of Open Mathematics Problems}, catalogue entry. CC BY 4.0. \url{https://huggingface.co/datasets/ulamai/UnsolvedMath}.
\item[\hypertarget{source:statement:2}{S2}] W. K. Hayman and E. F. Lingham, \emph{Research Problems in Function Theory} (2018), Chapter 6, Problem 6.59 and its complete update; Chapter 6 initial notation. \url{https://arxiv.org/abs/1809.07200}.
\end{description}
\label{end:statement}
\end{document}
